\section{Numerical details}
\label{app:numerics}

For complete simulation protocols, see Sec.~\ref{sec:protocols}.

\textit{Code.} The simulator stores 64 sites per 64-bit word. The eight neighbor bits of all sites
of a word are summed with bitwise adder trees, giving the number of live neighbors $k$ in four
bit planes; the flip class of each site is $(\sigma, k)$, with the flip probability $f_{\sigma k}$
read from an $2\times 9$ table, so that any outer-totalistic rule, Eq.~\eqref{M-eq:rule}, and the lazy
tables of Sec.~\ref{sec:async} are run by the same kernel. Each site draws one 32-bit uniform per
step from Philox$4\times32$-10~\cite{salmon2011} with counter (time, row, word, draw), and the 64
comparisons with the per-class thresholds are done bit-serially in parallel. Validation: (i) a
bit-exact comparison with a naive per-site implementation for five rules on four lattice shapes,
including one-word-wide and three-row lattices; (ii) per-class flip frequencies for an 18-class
table (variance of the $z$ scores over 560 repetitions $1.04\pm0.06$) and a $1.3\times10^8$-sample
test of the shared comparison on a voter-like table (per-class frequencies and adjacent-lane
independence, $|z|<2$); (iii) identical output for any number of threads and across restarts
from checkpoints; (iv) agreement with an independent GPU implementation of the $p = 0$ model to
1\% in the stationary density.

\textit{CUDA extension.} The dense GPU backend uses the original CPU integer
flip function and Philox implementation, extracted automatically into device
code. Two lattice buffers implement simultaneous updates; batches of 256 steps
are submitted as CUDA graphs without changing the update order or counters.
The host retains initialization, observables, profiles, and checkpoint format.
Validation compares full configurations, time series, profiles, snapshots, and
checkpoints byte for byte with the CPU implementation. Cases include one-word
rows, an $L=2048$ lattice, Moore and von Neumann neighborhoods, lazy update
tables, irregular output schedules, graph-batch boundaries, and counters near
$2^{32}$. CPU-to-GPU and GPU-to-CPU restarts reproduce the same trajectories.
The GPU checks absorption at scheduled outputs, so its absorption times are
interval observed. The lifetime and sparse-replica campaigns use the CPU code.
Parameter manifests record distinct 64-bit seeds, source and executable hashes,
and completion status; interrupted runs are identified rather than silently
pooled with completed ensembles.

\textit{Replica mode.} For seeds and droplets, only the bounding box of the sites with a
disagreeing neighbor (plus a margin) is updated, refitted every 32 steps. This is exact because a
site with $j = 0$ cannot flip and its random numbers are indexed by position, not by the order
of updates; over 200 seeds, box and full-lattice updating give identical configurations at
every step. A replica whose box comes within 2 rows or 64 columns of the lattice edge is stopped
and recorded as an escape. Its later survival is unobserved; the
seed-curve convention retains it as a separately tabulated contribution,
while the droplet growth observable counts the escape itself. In the pure voter model the method reproduces the martingale
$\langle A(t)\rangle = 4$ for the number of live sites of a $2\times2$ seed (within $|z|<0.94$
at all 38 output times over $2\times10^5$ replicas) and $P_s(t)\propto\ln t/t$.

\textit{Front tracking.} The broad strip survey in Table~\ref{tab:line}
uses a smoothed row profile and local interpolation of the velocities bracketing
reversal. For the orientation, extended-strip, and $p=1/2$ lazy-update
analyses, we explicitly track
the outer envelope: the lowest and highest bins whose smoothed activity exceeds
$0.3$ of the interior reference value measured at $t=10^3$. An interior hole
cannot stop this envelope. We also retain the center-connected-component width,
interior activity, and hole fraction as diagnostics. No run is removed simply
because its interior density is small in these corrected analyses.
The other supporting parameter surveys retain their legacy center-connected
tracker and interior-activity validity criterion; their crossing estimates
are not assigned the precision of the extended campaign.

For a tilted band, $s=(ay-bx)\bmod L$ has $L$ bins containing $L$ sites each,
with normal spacing $1/\sqrt{a^2+b^2}$. Smoothing is over nine normal lattice
spacings, converted to an odd number of bins. The front speed is half the slope
of the normal width. Fits stop when the band leaves the width range
$[0.15L,0.85L]$ in bin coordinates, before winding or disappearance invalidates
the geometry. In the broad orientation survey the principal window begins at
$t=2\times10^4$; fits starting at $10^4$ and $4\times10^4$ and threshold
fractions 0.2, 0.3, and 0.5 quantify sensitivity. Local bracketing estimates,
linear fits over wider intervals, and quadratic fits are compared explicitly.
Independent runs are resampled within each $q$ for the crossing intervals;
uncertainty from the fit model or time window is reported separately. The
extended campaign additionally compares late and nonoverlapping time windows;
each published fit requires full coverage of its declared window for every seed.

\textit{Sampling uncertainty.} Independent realizations are the statistical
units; lattice sites and successive output records are not treated as
independent samples. For some scalar means and paired mean differences in the
four-run ensembles, the reductions enumerate all $4^4=256$ ordered
empirical bootstrap resamples; other retained reductions use Monte Carlo
resampling with the count recorded in their provenance. Enumeration removes Monte Carlo error in the
resampling calculation, not the limited coverage accuracy from four runs.
Front-crossing intervals additionally condition on the stated interpolation,
window, and geometrical eligibility. Pointwise intervals are distinguished
from the simultaneous orientation contrasts. Exact binomial intervals are
used for fixed-count survival endpoints. All these sampling intervals remain
separate from finite-size, finite-time, and estimator sensitivity.

\textit{Spatial correlations.} The large-system correlations are evaluated
directly from the final packed-bit configurations, using periodic shifts in
each coordinate direction. Each stored configuration is checked against its
recorded size, time, digest, magnetization, and nearest-neighbor activity;
the identity $C(1)=1-2\rho$ supplies an additional consistency check.
All declared runs are retained. The same resampled run indices are used at
every separation when constructing the pointwise uncertainty bands.
These final-time measurements do not assume an exponential correlation form.

\textit{Droplets and holes.} For a row with live count $n_y$, the signed chord is $c_y=2n_y$ for an
active disk and $c_y=L-2n_y$ for an absorbing hole. The primary estimator is
$A=\sum_{|y-L/2|<1.3R_0+50}\max(c_y,0)$ with $R_0=1000$, followed by
$R=\sqrt{A/\pi}$. The crop reduces noise from the surrounding active phase;
rectification can introduce a positive area bias. We therefore also compute
$A_{\rm full}=\sum_y c_y$ and the trimmed signed sum, and report the resulting
rate sensitivity. Every such area stays positive in the retained trajectories.
Radius slopes are fitted over $2\times10^4\le t\le2\times10^5$.
The curvature overlays integrate the radius-dependent equation throughout
this window, with their initial radius fixed by the observed mean.


\textit{Layer edges.} For the roughness of the wetting layer (Sec.~\ref{sec:swall}) full
configurations were stored every $10^5$ steps. In each column the two edges of the layer are its
lowest and highest interface sites within $L_y/4$ of the wall; their variances along the wall and
the height-difference correlation $\langle[h(x+r)-h(x)]^2\rangle$ give the roughness.
These are two physical edges in the wall geometry. The same two-extrema
rule applied to a half-active band mixes its exterior front with an artificial
window boundary inside the active bulk. We therefore omit the previous
free-band comparator. The original band snapshots needed to replace that
curve are absent from the local data release; no corrected numerical
comparison is inferred from other snapshots.

\textit{Cost.} The earlier CPU campaigns used approximately 2200 core hours,
on Intel Xeon Platinum 8592+ nodes and a workstation, at
$(0.9$--$2.8)\times10^9$ site updates per second per core. The additional
campaign used 58.4 GPU allocation hours on NVIDIA H200 devices and
99.0 CPU-only allocated core hours. These allocation totals include validation,
I/O, and failed or restarted attempts; GPU-host CPU allocations are recorded
separately. Across the 40 staged GPU workers, mean sampled utilization during
simulation runs ranged from 92.7\% to 99.9\%. These samples include in-run
I/O but exclude build, validation, inter-task staging, and each run's initial
20 seconds; they are not allocation-wide efficiency estimates.

The additional front uncertainty check linearizes each locally interpolated
zero with respect to the two bracketing group means. It combines independent
within-group sample variances with Welch--Satterthwaite degrees of freedom.
For temporal contrasts, the same run's early and late linearized contributions
are paired before estimating variance. The angular contrast applies a
Bonferroni adjustment for the three pairwise direction comparisons. This is
a sensitivity calculation conditional on the observed bracket; it cannot
supply exact coverage from four runs or replace an assessment of systematic
front and finite-time effects.
